\documentclass[aspectratio=169,10pt,spanish]{beamer}

\input{../../../_preambulo_beamer.tex}
\input{../../../_comandos_beamer.tex}

\selectlanguage{spanish}

\title{MA1001B - Modelación Estadística para la Toma de Decisiones}
\subtitle{Lección 25: Distribución muestral de una proporción}
\author{}
\institute{Tecnol\'ogico de Monterrey}
\date{}

\begin{document}

\begin{frame}[plain]
  \vspace{1.2cm}
  {\Large\bfseries MA1001B - Modelación Estadística para la Toma de Decisiones\par}
  \vspace{0.7cm}
  {\large Lección 25: Distribución muestral de una proporción\par}
  \vspace{0.7cm}
  {\color{TecRojo}\rule{\textwidth}{1.2pt}}
  \vspace{0.8cm}

  Facultad MA1001B\\[0.3cm]
  Periodo\\[0.3cm]
  Tecnol\'ogico de Monterrey
\end{frame}

\begin{frame}{La pregunta de la proporción muestral}
  \begin{alertblock}{Pregunta guía}
    Si $p=0.18$ de los tickets llegan tarde, ¿cuánto se mueve $\hat p$ en
    una muestra de $n=120$, y cuándo es una campana la imagen equivocada?
  \end{alertblock}

  \vspace{0.6em}
  Al final de esta lección, usted puede:
  \begin{itemize}
    \item escribir $E[\hat p]=p$ y $\mathrm{EE}=\sqrt{p(1-p)/n}$;
    \item verificar las condiciones $np$ y $n(1-p)$;
    \item recuperar el EE con una simulación sembrada;
    \item explicar por qué $np<5$ rompe el modelo normal de $\hat p$.
  \end{itemize}

  \vfill
  \scriptsize Todos los indicadores de entrega tardía usados hoy son sintéticos. No se usan datos reales de empresa.
\end{frame}

\begin{frame}{EE de $\hat p$ y la verificación $np$}
  \small
  Población $p=0.18$, tamaño de muestra $n=120$:
  \[
    \mathrm{EE}(\hat p)=\sqrt{\frac{0.18\times 0.82}{120}}=0.035071.
  \]
  Conteos:
  \[
    np=21.6,\qquad n(1-p)=98.4.
  \]
  Ambos superan 5, de modo que una imagen normal de $\hat p$ es el modelo
  de trabajo.

  \vspace{0.3em}
  \begin{exampleblock}{Indicador sintético de entrega tardía}
    $\hat p$ es un estadístico muestral. $p$ es la tasa desconocida que se
    estudia.
  \end{exampleblock}
\end{frame}

\begin{frame}[fragile]{Paso 1: EE de fórmula}
  \begin{columns}[T]
    \begin{column}{0.42\textwidth}
      \small
      \begin{lstlisting}
se = sqrt(
  p * (1 - p) / n
)
np_count = n * p
      \end{lstlisting}

      \begin{block}{Salida verificada}
        $np=21.600000$\\
        $n(1-p)=98.400000$\\
        EE $=0.035071$
      \end{block}
    \end{column}
    \begin{column}{0.58\textwidth}
      \centering
      \includegraphics[width=\textwidth]{../figuras/sampling_proportion_01_se_phat.png}
      \vspace{0.2em}
      \scriptsize Modelo normal de $\hat p$ del Paso 1
    \end{column}
  \end{columns}
\end{frame}

\begin{frame}{La simulación recupera el EE}
  \small
  5{,}000 muestras con semilla 42 de $n=120$:
  \[
    \text{media de }\hat p=0.179663,\qquad
    \widehat{\mathrm{EE}}=0.035079.
  \]
  $P(\hat p>0.22)$ simulada $=0.119800$ versus normal $0.127032$.
\end{frame}

\begin{frame}[fragile]{Paso 2: 5{,}000 valores de $\hat p$}
  \begin{columns}[T]
    \begin{column}{0.40\textwidth}
      \small
      \begin{lstlisting}
successes = rng.binomial(
  n=120, p=0.18, size=5000
)
phat = successes / 120
      \end{lstlisting}

      \begin{block}{Salida verificada}
        Media $=0.179663$\\
        EE empírico $=0.035079$
      \end{block}
    \end{column}
    \begin{column}{0.60\textwidth}
      \centering
      \includegraphics[width=\textwidth]{../figuras/sampling_proportion_02_simulation.png}
      \vspace{0.2em}
      \scriptsize $\hat p$ simulada del Paso 2
    \end{column}
  \end{columns}
\end{frame}

\begin{frame}{Paso 3: $np<5$ rompe la campana}
  \begin{columns}[T]
    \begin{column}{0.42\textwidth}
      \small
      Ahora $n=20$, de modo que $np=3.6<5$.

      \vspace{0.4em}
      $P(\hat p=0)$ exacta $=0.018892$.\\
      Masa normal cerca de 0 $=0.035594$.\\
      Asimetría de $\hat p=0.430019$.

      \vspace{0.4em}
      \textbf{Límite:} $np<5$ hace fallar la aproximación normal.
      La Lección 26 estudia estimadores puntuales.
    \end{column}
    \begin{column}{0.58\textwidth}
      \centering
      \includegraphics[width=\textwidth]{../figuras/sampling_proportion_03_np_condition_limit.png}
      \vspace{0.2em}
      \scriptsize $\hat p$ discreta versus campana ilegal del Paso 3
    \end{column}
  \end{columns}
\end{frame}

\begin{frame}{Verificación de conceptos}
  \begin{enumerate}
    \item Calcule $\mathrm{EE}(\hat p)$ para $p=0.18$ y $n=120$.
    \item ¿Por qué importan $np=21.6$ y $n(1-p)=98.4$?
    \item El EE empírico es $0.035079$. ¿Sustituye eso a $0.035071$?
    \item ¿Por qué $P(\hat p=0)=0.018892$ es un problema para un modelo
      normal?
  \end{enumerate}

  \vspace{0.6em}
  \begin{block}{Conexión con la Lección 26}
    La sesión puede continuar directamente con insesgadez, eficiencia y
    consistencia de estimadores puntuales.
  \end{block}

  \vfill
  \scriptsize Fuente: OpenStax, \textit{Introductory Business Statistics 2e},
  Capítulo 7, Sección 7.3. CC BY-NC-SA.
\end{frame}

\appendix



\end{document}
